% Working outline adapted from Ripple's paper structure.
% Choose the venue and replace prompts after the experiments establish results.
\documentclass[conference]{IEEEtran}
\usepackage{times}
\usepackage[numbers]{natbib}
\usepackage{booktabs}
\usepackage{tabularx}
\usepackage[bookmarks=true]{hyperref}

\newcommand{\stub}[1]{\par\noindent\textit{TODO: #1}\par}

\begin{document}

\title{SCRIMMAGE-RS: A Rust Testbed for Reproducible Multi-Agent Robotics Experiments}
\author{Author names and affiliations to be added}
\maketitle

\begin{abstract}
\stub{State the research problem, the retained SCRIMMAGE design, and the
changes evaluated in Rust. Summarize the C++ comparison, one demonstrated
research extension, and measured outcomes after the experiments are complete.}
\end{abstract}

\IEEEpeerreviewmaketitle

\section{Introduction}

SCRIMMAGE provides a plugin-based environment for multi-agent robotics
experiments \cite{demarco2018}. This paper will evaluate a Rust implementation
that retains its useful simulation structure while simplifying mission
configuration and the path from a user plugin to repeatable experiments.

\stub{Introduce the research use case and concrete difficulties motivating the
changes. Explain what users gain from a single mission description, workspace
plugins, and shared run, sweep, and Slurm tooling.}

The planned evaluation asks:
\begin{description}
  \item[RQ1:] Where do the C++ and Rust implementations agree in behavior, and
  how do their runtime and memory costs change with workload and agent count?
  \item[RQ2:] What work does the Rust configuration and plugin workflow remove
  when building and repeating a multi-agent experiment?
  \item[RQ3:] What additional experiment does a selected new capability enable,
  and what does that capability cost?
\end{description}

\stub{State the final contributions once supported by evidence.}

\section{How the Field Reached This Point}
\label{sec:background}

\subsection{The Original SCRIMMAGE Design}
\stub{Review the original paper's design guidelines, plugin architecture,
integrations, and predator/prey study. Distinguish the published system from
the current C++ branch used for the comparison.}

\subsection{Current Research Needs and Related Systems}
\stub{Select recent primary sources relevant to the chosen experiment, such as
learned or hybrid autonomy, batched inference, and reproducible multi-agent
evaluation. Compare their assumptions and workloads. Add citations as sources
are reviewed; library availability alone does not establish a contribution.}

\section{SCRIMMAGE-RS}
\label{sec:system}

\subsection{Experiment and Data Path}
\stub{Show one small YAML mission with templates, a parameter sweep, and the
path through validation, execution, saved results, and Rerun inspection.
Explain how XML and YAML share construction. Add one dataflow figure.}

\subsection{Plugin Authoring and Execution}
\stub{Describe the seven model roles, workspace starter, compiled registration,
entity ownership, phase barriers, and message timing. Use a small plugin
example to show what the researcher writes and what the engine handles.}

\subsection{Shared Research Tools}
\stub{Walk through the same custom plugin running locally and in a Slurm
campaign through the shared command and tools. Show how framework upgrades
reach workspace experiments without copying runners or schedulers.}

\subsection{Selected Research Extension}
\stub{Choose one concrete experiment. Candidates from the roadmap include
Burn-backed policy inference, a ROS or ArduPilot connection, or direct typed
Rust construction. These remain planned work. Describe the implemented
boundary, research use, and comparison only when a pilot exists.}

\subsection{Validation}
\stub{Explain analytical and plugin tests, selected C++ trajectory and summary
comparisons, and exact Rust frame/event/summary checks at 1, 2, and 8 workers.
Describe XML/YAML equivalence and recording on/off checks. List known
differences and the absence of direct C++ event-stream comparison.}

\section{Benchmark Design}
\label{sec:benchmark}

\subsection{Baseline and Workloads}
\stub{Select a small shared set: simple motion, controller substeps, sensing or
communication, and spawn/removal. Scale populations and model complexity
separately. Record equivalent parameters, initial states, timestep, seeds,
horizon, plugin behavior, and output settings for both implementations.}

\subsection{Measurements}
\stub{Measure end-to-end wall time, startup, simulation cost where instrumented,
peak memory, output size, and failures. Report simulated seconds per wall
second and worker scaling. Treat unsupported or failed cases explicitly.}

\subsection{Configuration and Extension Effort}
\stub{Use the same concrete tasks in both systems: add an agent type, change a
population, add a custom plugin, and launch a sweep. Show files and edits
required. Keep descriptive code comparisons separate from claims about human
usability unless a user study is performed.}

\subsection{New Capability Evaluation}
\stub{Compare the selected extension against a useful baseline and isolate its
effect. For accelerated inference, include packing, transfers, synchronization,
and readback in full-run cost; for an external controller, report timing and
failure behavior as well as task performance.}

\section{Experimental Design}
\label{sec:experiments}

\stub{Freeze the C++ branch revision and Rust revision for the campaign. Record
hardware, OS, compiler versions, release flags, dependencies, threading,
container architecture, and whether execution is native or emulated. Run both
implementations on the same hardware and execution environment for timing.}

\stub{Check behavior before interpreting speed differences. Save initial states
when random generators differ. Separate deliberate semantic differences from
regressions, and report the covered subset of models.}

\stub{Specify workload sizes, independent seeds, warm-up policy, timing repeats,
run ordering, and uncertainty estimates. Separate repeated timing measurements
of one case from independent task trials. Compare recording modes explicitly.}

\stub{Use the existing reference and bulk-run tools. Save commands, inputs,
raw measurements, and failures. Complete an actual Slurm campaign before
claiming cluster validation.}

\section{Results}
\label{sec:results}

\subsection{Behavioral Agreement and Differences}
\stub{Add the comparison matrix, tolerances, failures, and explanations.}

\subsection{Runtime, Memory, and Scaling}
\stub{Add measured curves and uncertainty. Separate one-worker C++/Rust
comparisons from Rust worker scaling and from recording overhead.}

\subsection{Configuration and Experiment Workflow}
\stub{Show the worked configuration/plugin comparison and campaign results.}

\subsection{Selected Research Extension}
\stub{Report task outcomes, integration effort, and complete execution cost.}

\section{Discussion and Limitations}

\stub{Explain what the measurements support and where they do not generalize:
model coverage, fidelity, stochastic differences, platform effects, bottlenecks,
and external integrations. Distinguish gains due to implementation choices
from claims about Rust as a language. Compare benefits with added complexity.}

\section{Conclusion}

\stub{Answer the research questions using measured results and identify the
most useful next experiment.}

\section*{Generative AI Disclosure}

OpenAI Codex assisted with the initial manuscript outline and reference-file
organization. \stub{Update this statement to reflect actual assistance and
author review, following the selected venue's policy.}

\bibliographystyle{plainnat}
\bibliography{../references}

\end{document}
